\documentclass[11pt]{article}

\usepackage[T1]{fontenc}
\usepackage[margin=1in]{geometry}
\usepackage{amsmath,amssymb,amsthm}
\usepackage[numbers]{natbib}
\usepackage[hidelinks]{hyperref}
\usepackage{microtype}

\newtheorem{theorem}{Theorem}[section]
\newtheorem{proposition}[theorem]{Proposition}
\newtheorem{lemma}[theorem]{Lemma}
\newtheorem{corollary}[theorem]{Corollary}
\theoremstyle{definition}
\newtheorem{definition}[theorem]{Definition}
\theoremstyle{remark}
\newtheorem{remark}[theorem]{Remark}

\newcommand{\F}{\mathbb F}
\newcommand{\AS}{\operatorname{AS}}
\newcommand{\Tr}{\operatorname{Tr}}
\newcommand{\Push}{\operatorname{Push}}
\newcommand{\Lift}{\operatorname{Lift}}

\title{Fast Multiplication in Canonical Nim Coordinates}
\author{GPT 5.6 Sol \and a9lim}
\date{}

\begin{document}
\maketitle

\begin{abstract}
For $n=2^k$, the finite nimbers below $2^n$ form $\F_{2^n}$, but their
literal $n$-bit words are not an arbitrary polynomial basis.  They are the
multivariate basis of the Conway Artin--Schreier tower generated by
$c_i=2^{2^i}$ with
\[
  c_i^2+c_i=\prod_{j<i}c_j.
\]
We give explicit forward and inverse transforms from these canonical words to
the primitive Artin--Schreier tower of De Feo and Schost.  The specialization
has a particularly simple form: the image of each canonical generator is an
affine shift $z_i+r_i$ of the primitive generator, and the two coefficient
blocks change by
\[
  (A,B)\longmapsto(A+Br_i,B).
\]
The shift is its own inverse in characteristic two.

After uniform preprocessing, the resulting multiplication algorithm uses
\[
  O\bigl(M(n)\log n+n\log^3 n\bigr)
\]
bit operations and $n\log^{O(1)}n$ space, where $M(n)$ is binary-polynomial
multiplication cost.  Precomputation is also quasi-linear up to logarithmic
factors and stores no dense change-of-basis matrix or multiplication tensor.
The inverse transform returns the original tower coefficients, whose ordering
is exactly the literal nim-word ordering; the output is therefore the
canonical nim product, not merely an element of an abstract isomorphic field.

This closes the mathematical transform and complexity part of the problem.
An arbitrary-width implementation, differential tests, and crossover
measurements remain engineering work.  Lean checks the affine
Artin--Schreier identity, the involutive block transform, preservation of the
tower equations, and exact multiplication transport; the asymptotic
algorithms and the concrete Conway identification remain paper-level and
cited mathematics.
\end{abstract}

\section{Statement of the result}

Nim addition and multiplication make the ordinals a characteristic-two
field-like proper class; the finite segment below $2^{2^k}$ is the field
$\F_{2^{2^k}}$ \cite{conway1976,lenstra1978}.  Put
\[
  n=2^k,
  \qquad
  K_k=\{0,\ldots,2^n-1\}=\F_{2^n}.
\]
An element of $K_k$ is supplied as its literal $n$-bit integer word.  The
problem is to multiply in time quasi-linear in this representation length,
without replacing the canonical word by an unspecified field presentation or
hiding a quadratic table in preprocessing.

We use a standard multiplication function $M$ satisfying
\[
  M(m)\geq m,
  \qquad
  \frac{M(m)}m\text{ is nondecreasing},
  \qquad
  M(2m)=O(M(m)).
\]
We also choose the usual fast polynomial multiplication, remainder, and
transposed algorithms with $m\log^{O(1)}m$ workspace.  These are the same
regularity assumptions used in the fast-tower analysis
\cite{defeoschost2012}.

\begin{theorem}[canonical-word multiplication]\label{thm:main}
For every $k\geq0$, let $n=2^k$.  There is a uniform preprocessing algorithm
which constructs $n\log^{O(1)}n$ bits of level data, followed by a
multiplication algorithm
\[
  \operatorname{mul}_k:K_k\times K_k\longrightarrow K_k
\]
on literal canonical nim words.  Its online cost is
\[
  O\bigl(kM(n)+nk^3\bigr)
  =O\bigl(M(n)\log^{O(1)}n\bigr),
\]
and its auxiliary space is $n\log^{O(1)}n$ bits.  The preprocessing cost is
\[
  O\bigl(k^2M(n)+nk^4+M(2n)\bigr).
\]
No dense basis-conversion matrix or multiplication tensor is stored.
\end{theorem}

The proof has three parts.  First, the literal bit ordering is identified with
the monomial ordering of one explicit quadratic tower.  Second, that tower is
mapped to the primitive tower of De Feo and Schost by recursively computed
affine generator shifts.  Third, the published level-embedding bounds are
specialized and summed without suppressing conversion or preprocessing.

\section{The literal word is the Conway tower basis}

It is useful to index generators from one.  Set
\[
  u_i=c_{i-1}=2^{2^{i-1}}
  \quad (i\geq1),
  \qquad
  A_1=1,
  \qquad
  A_i=\prod_{j<i}u_j\quad(i\geq2).
\]
The canonical finite-nimber identities are
\begin{equation}
  u_i^2+u_i=A_i.
\end{equation}
Equivalently, $c_i^2=c_i+\prod_{j<i}c_j$.  This is the finite
Conway--Lenstra tower relation, not an arbitrary choice of irreducible
polynomials \cite{conway1976,lenstra1978}.

Define the presented tower
\[
  T_0=\F_2,
  \qquad
  T_i=T_{i-1}[U_i]/(U_i^2+U_i+A_i),
\]
where after adjoining $u_i$ we put $A_{i+1}=A_i u_i$.

\begin{proposition}[trace-one sources]\label{prop:trace-one}
For every $i\geq1$,
\[
  \Tr_{T_{i-1}/\F_2}(A_i)=1.
\]
Consequently $X^2+X+A_i$ is irreducible over $T_{i-1}$,
$[T_i:T_{i-1}]=2$, and the $2^i$ monomials
\[
  \prod_{j=1}^i u_j^{\epsilon_j},
  \qquad \epsilon_j\in\{0,1\},
\]
form an $\F_2$-basis of $T_i$.
\end{proposition}

\begin{proof}
We prove simultaneously that every $T_i$ is a field and that the next source
has trace one.  The initial field is $T_0=\F_2$ and $A_1=1$.  Suppose
$T_{i-1}$ is a field and the trace of $A_i$ is one.  The finite-field
Artin--Schreier criterion makes $X^2+X+A_i$ irreducible, so $T_i$ is a
quadratic field extension.  Its two conjugates of $u_i$ are $u_i$ and
$u_i+1$, whence
\[
  \Tr_{T_i/T_{i-1}}(u_i)=1.
\]
It follows that
\[
  \Tr_{T_i/T_{i-1}}(A_{i+1})
  =\Tr_{T_i/T_{i-1}}(A_i u_i)=A_i.
\]
Trace transitivity proves the displayed trace identity at level $i+1$.
This completes the induction by trace transitivity.  The criterion used above
is the standard equivalence between irreducibility of $X^2+X+a$ and absolute
trace one \cite{lidlniederreiter1997}.  Every step therefore doubles the
degree, and the stated monomials are the resulting iterated quadratic basis.
\end{proof}

\begin{proposition}[literal ordering]\label{prop:literal-order}
Let $0\leq e<2^k$ and write
\[
  e=\sum_{i=1}^k \epsilon_i2^{i-1},
  \qquad \epsilon_i\in\{0,1\}.
\]
Then the nimber represented by the lone word bit in position $e$ is
\[
  2^e=\prod_{i=1}^k u_i^{\epsilon_i}
\]
under nim multiplication.  Hence, for a word
$x=\sum_{e<n}x_e2^e$, the bit vector $(x_e)$ is exactly the coordinate
vector of $x$ on the monomial basis of Proposition~\ref{prop:trace-one}, in
ordinary increasing bit order.
\end{proposition}

\begin{proof}
Distinct Fermat powers nim-multiply as their ordinary product, while nim
addition of finite nimbers is bitwise exclusive-or
\cite{conway1976,lenstra1978}.  Since $u_i=2^{2^{i-1}}$, the displayed
product is $2^{\sum_i\epsilon_i2^{i-1}}=2^e$.  Distributing over the set
bits of $x$ gives the coordinate statement.
\end{proof}

This proposition is the representation boundary that an abstract
finite-field isomorphism would not supply.  From now on, splitting an
$N_i=2^i$ bit canonical word into its low and high halves means writing
\begin{equation}
  x=x_0+x_1u_i,
  \qquad x_0,x_1\in T_{i-1}.
\end{equation}

\section{The fast primitive tower}

We specialize the primitive Artin--Schreier tower of De Feo and Schost to
$p=2$, base degree $d=1$, and $Q_0=X+1$
\cite{defeoschost2012}.  Let $V_0=\F_2$ and
\[
  V_i=V_{i-1}(z_i),
  \qquad
  z_i^2+z_i=g_i,
\]
where
\begin{equation}
  g_1=1,
  \qquad
  g_2=z_1,
  \qquad
  g_i=z_{i-1}^3\quad(i\geq3).
\end{equation}
Their primitive-tower theorem says that $z_i$ generates all of $V_i$ over
$\F_2$.  If $Q_i$ is its minimal polynomial, then
\[
  V_i=\F_2[Z]/(Q_i(Z))
\]
with the univariate power basis
\[
  1,z_i,\ldots,z_i^{N_i-1}.
\]
Multiplication and reduction in this basis cost $O(M(N_i))$ bit operations.

The same work gives mutually inverse level maps
\[
  \Push_i:V_i\longrightarrow V_{i-1}\oplus V_{i-1}z_i,
  \qquad
  \Lift_i=\Push_i^{-1},
\]
with cost
\begin{equation}
  L(i)=O\bigl(N_i i^2+M(N_i)\bigr).
\end{equation}
They are implemented by divide-and-conquer reduction, transposed reduction,
and trace parameterization; they are algorithms, not stored matrices.  The
construction also computes the $Q_i$, modular-reduction data, traces, and
pseudotraces used by the level maps and the fast Artin--Schreier solver.

The primitive tower is not the canonical nim tower.  The first two sources in
the canonical and primitive displays
coincide, but at the third step the primitive source is $z_2^3$, whereas the
canonical source is the product of the first two generators.  The conversion
below is therefore essential.

\section{Affine-shift isomorphisms}

Write
\[
  \AS(x)=x^2+x.
\]
In characteristic two,
\begin{equation}
  \AS(x+y)=\AS(x)+\AS(y).
\end{equation}
We now construct compatible field isomorphisms
\[
  \sigma_i:T_i\longrightarrow V_i
\]
and expose their exact coordinate action.

\begin{theorem}[selected affine shifts]\label{thm:shifts}
Set $\sigma_0=\operatorname{id}_{\F_2}$.  Suppose $\sigma_{i-1}$ has been
constructed, and put
\[
  S_i=\sigma_{i-1}(A_i)
      =\prod_{j<i}s_j,
  \qquad
  s_j=\sigma_j(u_j).
\]
There is an $r_i\in V_{i-1}$ satisfying
\begin{equation}
  r_i^2+r_i=S_i+g_i.
\end{equation}
For any such $r_i$, the assignment
\begin{equation}
  s_i=z_i+r_i
\end{equation}
extends $\sigma_{i-1}$ uniquely to a field isomorphism $\sigma_i:T_i\to V_i$.
The shifts $r_i$ are computable uniformly by the fast Artin--Schreier solver.
\end{theorem}

\begin{proof}
Proposition~\ref{prop:trace-one} and invariance of trace under field
isomorphism give
\[
  \Tr_{V_{i-1}/\F_2}(S_i)=1.
\]
The primitive-tower theorem gives the same identity for $g_i$.  Therefore
$S_i+g_i$ has trace zero.  The image of $x\mapsto x^2+x$ on a finite
characteristic-two field is exactly the absolute-trace kernel, so
the displayed shift equation has a solution.  De Feo and Schost give a
recursive quasi-linear algorithm for finding one.

The primitive relation, shift equation, and additivity of $\AS$ yield
\[
  s_i^2+s_i
  =g_i+(S_i+g_i)=S_i.
\]
Thus $s_i$ satisfies the transported defining equation of $u_i$.  The quotient
universal property extends $\sigma_{i-1}$ to a homomorphism $T_i\to V_i$.
Both fields have $2^{N_i}$ elements, and a unital field homomorphism is
injective, so it is an isomorphism.
\end{proof}

The affine form is stronger than merely knowing some coordinates for $s_i$.
Indeed, any root in $V_i$ of an Artin--Schreier polynomial with source in
$V_{i-1}$ is either downstairs or has $z_i$-coefficient one.  The selected
root is not downstairs because its source has trace one there.  Hence
the affine formula $s_i=z_i+r_i$ loses no generality and needs no stored
coefficient inverse.

\begin{remark}[first nontrivial shift]
Let $a=z_1$, $b=z_2$, and $t=z_3$.  Then
\[
  a^2+a=1,
  \qquad
  b^2+b=a,
  \qquad
  t^2+t=b^3=ab+a+b.
\]
Put $r_3=ab+a+b$.  Direct reduction gives $r_3^2+r_3=a+b$, and hence
\[
  (t+r_3)^2+(t+r_3)=ab,
\]
the canonical third source.  This is where the two towers first diverge.
\end{remark}

\section{Forward and inverse transforms}

For $A,B\in V_{i-1}$, the entire level conversion rests on
\begin{equation}
  A+B(z_i+r_i)=(A+Br_i)+Bz_i.
\end{equation}
The coefficient map
\begin{equation}
  \tau_i(A,B)=(A+Br_i,B)
\end{equation}
is an involution, since
\[
  (A+Br_i)+Br_i=A.
\]

Define the forward transform $F_i$ recursively.  Set $F_0$ to the identity.
Given the canonical low--high split, compute
\begin{align}
  X_0&=F_{i-1}(x_0),&
  X_1&=F_{i-1}(x_1),\notag\\
  F_i(x)&=\Lift_i(X_0+r_iX_1,X_1).
\end{align}
Thus $F_i(x)$ is returned in the univariate power basis of $z_i$.

Define the inverse $G_i$ in the opposite direction.  Set $G_0$ to the
identity.  For $y\in V_i$, compute
\[
  \Push_i(y)=(Y_0,Y_1),
\]
and put
\begin{equation}
  G_i(y)=G_{i-1}(Y_0+r_iY_1)+u_iG_{i-1}(Y_1).
\end{equation}
The right side is emitted as the low and high halves of a canonical word; it
does not require a tower multiplication.

\begin{proposition}[transform correctness]\label{prop:transform-correct}
For every $i$, $F_i$ is the coordinate algorithm for $\sigma_i$, $G_i$ is
the coordinate algorithm for $\sigma_i^{-1}$, and
\[
  G_iF_i=\operatorname{id}_{T_i},
  \qquad
  F_iG_i=\operatorname{id}_{V_i}.
\]
\end{proposition}

\begin{proof}
Assume the statement at level $i-1$.  Replacing $u_i$ by
$s_i=z_i+r_i$ in $x_0+x_1u_i$ gives the forward formula by the block-shift
identity, so $F_i$ computes $\sigma_i$.  Conversely, the involutive block map
rewrites the pushed-down coefficients of $y$
from the $z_i$ basis to the $s_i$ basis, after which
the inverse formula applies the lower inverse transforms.  Since
$\tau_i$ is involutive and the level maps $\Push_i,\Lift_i$ are inverse,
both displayed identities follow by induction.
\end{proof}

We may now define
\begin{equation}
  \operatorname{mul}_k(x,y)
  =G_k\bigl(F_k(x)F_k(y)\bmod Q_k\bigr).
\end{equation}

\begin{corollary}[exact canonical output]\label{cor:exact-output}
The bit string returned by the displayed multiplication algorithm is the literal
canonical nim product of $x$ and $y$.
\end{corollary}

\begin{proof}
Proposition~\ref{prop:transform-correct} says that $F_k$ and $G_k$ are inverse
field isomorphisms, so decoding the returned vector in $T_k$ gives the field
product of the decoded inputs.  Proposition~\ref{prop:literal-order} says that
the $T_k$ basis coordinates, in the order returned by $G_k$, are exactly the
literal nim-word bits.  Uniqueness of coordinates in that basis proves the
claim.
\end{proof}

\section{Time and space bounds}

Let $T(i)$ be the cost of either transform at level $i$.  A depth-first
implementation makes two recursive calls, one level conversion, and at most
one multiplication in $V_{i-1}$.  By the level-embedding bound,
\begin{equation}
  T(i)\leq2T(i-1)
    +O\bigl(L(i)+M(N_{i-1})\bigr).
\end{equation}
Unrolling at $n=N_k$ gives
\[
  T(k)=
  O\left(
    \sum_{i=1}^k2^{k-i}
      \bigl(N_i i^2+M(N_i)\bigr)
  \right).
\]
The first summand at depth $i$ is $ni^2$, so its total is $O(nk^3)$.
Monotonicity of $M(m)/m$ gives
\[
  2^{k-i}M(2^i)\leq M(2^k)=M(n),
\]
so the multiplication contribution is $O(kM(n))$.  Hence
\begin{equation}
  T(k)=O\bigl(nk^3+kM(n)\bigr).
\end{equation}
Two forward transforms, one $O(M(n))$ univariate multiplication, and one
inverse transform have the same asymptotic bound.  This proves the online
part of Theorem~\ref{thm:main}.

For completeness, the preprocessing is not free.  In the notation of
De Feo and Schost, the level-embedding cost is $L(i)$ and their pseudotrace
bound specializes at $p=2,d=1$ to
\[
  \operatorname{PT}(i)=O(i^2L(i)).
\]
Their Artin--Schreier solver has the same dominant bound.  Constructing the
primitive tower, its remainder and trace data, and all shifts through level
$k$ therefore costs
\[
  O\bigl(k^2L(k)+M(2n)\bigr)
  =O\bigl(k^2M(n)+nk^4+M(2n)\bigr).
\]
This is quasi-linear up to logarithms.  In particular, no quadratic-time
preprocessing is being exchanged for fast online multiplication.

The stored primitive polynomials, reciprocal data, trace series, selected
shifts, and level maps have total size $n\log^{O(1)}n$.  More explicitly, the
fast-tower construction stores $O(k^2)$ tower elements; charging every one the
top-level width gives the conservative bound $O(nk^2)$ bits.  The actual
level-weighted sum is smaller for data kept only at their construction level.
Push-down and lift-up are divide-and-conquer algorithms, not $N_i$-square
matrices.  Running the two recursive children sequentially retains only a
geometric stack of coefficient blocks, and the selected polynomial routines
use quasi-linear workspace.  Thus both working storage and conversion data
are $n\log^{O(1)}n$.

\section{Formal and implementation boundary}

The import-only paper entry point combines the shared
\texttt{Algebra.ArtinSchreier} module with the paper-specific
\texttt{NimFastMultiplication} module.  Together they check the following
algebraic joins used above:
\begin{itemize}
\item the finite-field equivalence between trace one and irreducibility of
  $X^2+X+a$;
\item the trace carry through a quadratic generator and the direct
  three-variable-product quadratic multiplication split;
\item the shifted-source identity of Theorem~\ref{thm:shifts};
\item the involutivity and evaluation identity for
  $(A,B)\mapsto(A+Br,B)$;
\item preservation of the canonical tower equations under a ring
  equivalence; and
\item exact recovery of basis coordinates after multiplication is transported
  through a ring equivalence.
\end{itemize}
The existing finite-field development also identifies the Artin--Schreier
image with the trace kernel.  These declarations are kernel-checked by Lean
and Mathlib \cite{demouraullrich2021,mathlib2020}.

Lean does not formalize the concrete class of finite nimbers, the
Conway--Lenstra identities, the De Feo--Schost divide-and-conquer algorithms,
or their asymptotic analysis.  Those are respectively paper-level cited
mathematics and external algorithmic theorems.  The formal module therefore
checks the load-bearing algebraic specialization without pretending to be an
end-to-end verification of Theorem~\ref{thm:main}.

Theorem~\ref{thm:main} and Corollary~\ref{cor:exact-output} supply the forward
and inverse transforms, their correctness, multiplication, preprocessing,
time, and space bounds.  The mathematical portion of fast multiplication in
canonical nim coordinates is consequently closed.  What remains is to
implement arbitrary-width bit vectors and polynomial arithmetic, compare
exhaustively with the mex oracle on small fields, compare differentially with
the current fixed-width backend, and measure the crossover at which the fast
path should replace it.

\clearpage
\small
\bibliographystyle{plainnat}
\bibliography{nim_fast_multiplication}

\end{document}
